% =====================================================================
% preamble.tex - Diseño de la versión PDF (A4) de R-Stars: El Universo.
% Solo lo usa bookdown::pdf_book (ver _output.yml). La web no lo lee.
% =====================================================================

% --- Idioma ----------------------------------------------------------
% El español lo activa pandoc (lang=es en _output.yml, mediante babel).
% Aquí solo se cambia "Cuadro" por "Tabla", como en la web.
\addto\captionsspanish{%
  \renewcommand{\tablename}{Tabla}%
  \renewcommand{\listtablename}{Índice de tablas}}

% --- Tipografía (la misma familia que la web) -------------------------
\usepackage{sourceserifpro}
\usepackage{sourcesanspro}
\usepackage[scale=0.86]{sourcecodepro}
\linespread{1.08}

% --- Página A4 --------------------------------------------------------
% Doble cara: los capítulos empiezan siempre en página impar (openright,
% en _output.yml). Márgenes interior y exterior iguales: vale también
% para imprimir a una cara.
\usepackage[a4paper, top=2.6cm, bottom=2.6cm, inner=2.5cm, outer=2.5cm,
            headheight=15pt, headsep=0.8cm]{geometry}
\usepackage{emptypage}   % las páginas que se dejan en blanco, sin cabecera
\setlength{\emergencystretch}{3em}   % menos líneas que se salen del margen
\raggedbottom
\widowpenalty=10000
\clubpenalty=10000

% --- Colores (los de la web: primary #475569) --------------------------
\usepackage{xcolor}
\definecolor{rstars}{HTML}{475569}
\definecolor{rstarsclaro}{HTML}{F1F5F9}
\definecolor{rstarsgris}{HTML}{64748B}
\definecolor{rstarsfila}{HTML}{F1F5F9}  % filas alternas de las tablas

% --- Títulos ----------------------------------------------------------
\usepackage{titlesec}
% Los títulos van sin justificar (alineados a la izquierda) y sin partir
% palabras con guion (p. ej. "MATrstars" no se rompe nunca).
\newcommand{\titulorstars}{\raggedright\hyphenpenalty=10000\exhyphenpenalty=10000}
\titleformat{\chapter}[display]
  {\sffamily\bfseries\color{rstars}\titulorstars}
  {\large\MakeUppercase{\chaptertitlename}\ \thechapter}
  {6pt}
  {\huge}
  [\vspace{4pt}{\color{rstars}\titlerule[1pt]}]
\titleformat{name=\chapter, numberless}[display]
  {\sffamily\bfseries\color{rstars}\titulorstars}{}{0pt}{\huge}
  [\vspace{4pt}{\color{rstars}\titlerule[1pt]}]
\titlespacing*{\chapter}{0pt}{0pt}{28pt}
\titleformat{\section}
  {\sffamily\Large\bfseries\color{rstars}\titulorstars}{\thesection}{0.7em}{}
\titleformat{\subsection}
  {\sffamily\large\bfseries\color{rstars}\titulorstars}{\thesubsection}{0.7em}{}
% Un título nunca se queda solo al pie de una página: si no caben tras él
% unas líneas de su contenido, pasa a la página siguiente.
\usepackage{needspace}
\usepackage{etoolbox}
\preto\section{\Needspace{9\baselineskip}}
\preto\subsection{\Needspace{7\baselineskip}}
\preto\subsubsection{\Needspace{6\baselineskip}}
\titleformat{\subsubsection}
  {\sffamily\normalsize\bfseries\titulorstars}{\thesubsubsection}{0.7em}{}

% --- Cabeceras y pies -------------------------------------------------
\usepackage{fancyhdr}
\pagestyle{fancy}
\fancyhf{}
\fancyhead[LE]{\small\sffamily\color{rstarsgris}R-Stars: El Universo}
\fancyhead[RO]{\small\sffamily\color{rstarsgris}\nouppercase{\leftmark}}
\fancyfoot[C]{\small\sffamily\thepage}
\renewcommand{\headrulewidth}{0.4pt}
\renewcommand{\chaptermark}[1]{\markboth{#1}{}}
\fancypagestyle{plain}{%
  \fancyhf{}\fancyfoot[C]{\small\sffamily\thepage}%
  \renewcommand{\headrulewidth}{0pt}}

% --- Figuras y tablas: en su sitio, junto a su rótulo -------------------
\usepackage{float}
\floatplacement{figure}{H}
\floatplacement{table}{H}

% Figura o tabla + su rótulo: pdf/filtro.lua los encierra en un bloque
% que no se puede partir entre páginas. Si el bloque es más alto que la
% página (p. ej. una tabla muy larga), se reduce lo justo para que quepa.
\usepackage{calc}   % anchos de columna de las tablas de markdown
\usepackage{booktabs}
\usepackage{colortbl}
\newsavebox{\rstarsbloquecaja}
\newenvironment{rstarsbloque}
  {\par\medskip\begin{lrbox}{\rstarsbloquecaja}\begin{minipage}{\linewidth}\centering}
  {\end{minipage}\end{lrbox}%
   \typeout{RStars bloque: alto \the\dimexpr\ht\rstarsbloquecaja+\dp\rstarsbloquecaja\relax\space (pagina \the\textheight), pag. \thepage}%
   {\centering
    \ifdim\dimexpr\ht\rstarsbloquecaja+\dp\rstarsbloquecaja\relax>0.97\textheight
      \resizebox*{!}{0.97\textheight}{\usebox{\rstarsbloquecaja}}%
    \else
      \usebox{\rstarsbloquecaja}%
    \fi\par}%
   \medskip}

% Rótulos "Tabla X.Y." / "Figura X.Y." (los centra pdf/filtro.lua)
\newenvironment{rotulo}
  {\par\nopagebreak[4]\vspace{0.4\baselineskip}\begingroup\centering\small}
  {\par\endgroup\medskip}

% --- Código -----------------------------------------------------------
\usepackage{fvextra}
% Salida de consola de R: parte las líneas largas en vez de desbordar.
\DefineVerbatimEnvironment{SalidaR}{Verbatim}{%
  breaklines, breakanywhere, fontsize=\small,
  xleftmargin=0.8em, formatcom=\color{rstars}}
% Código R resaltado: mismo comportamiento de corte de línea.
\makeatletter
\AtBeginDocument{%
  \@ifundefined{Highlighting}{}{%
    \RecustomVerbatimEnvironment{Highlighting}{Verbatim}{%
      commandchars=\\\{\}, breaklines, breakanywhere, fontsize=\small}}}
\makeatother

% --- Tablas anchas -----------------------------------------------------
% pdf/filtro.lua envuelve cada tabla en \tablarstars{...}, que la
% mide: si cabe, la deja igual; si se pasa hasta un 40 %, la reduce al
% ancho del texto; si es más ancha, la gira 90° y la ajusta al alto de la
% página. El ancho de cada tabla queda anotado en el .log ("RStars tabla").
\newsavebox{\rstarscaja}
\newcommand{\tablarstars}[1]{%
  \sbox{\rstarscaja}{#1}%
  \typeout{RStars tabla: \the\wd\rstarscaja\space (texto \the\linewidth), pag. \thepage}%
  \ifdim\wd\rstarscaja>1.4\linewidth
    % muy ancha: se gira 90 grados y se ajusta al alto y al ancho disponibles
    \sbox{\rstarscaja}{\rotatebox{90}{\usebox{\rstarscaja}}}%
    \ifdim\dimexpr\ht\rstarscaja+\dp\rstarscaja\relax>0.85\textheight
      \sbox{\rstarscaja}{\resizebox*{!}{0.85\textheight}{\usebox{\rstarscaja}}}%
    \fi
    \ifdim\wd\rstarscaja>\linewidth
      \sbox{\rstarscaja}{\resizebox{\linewidth}{!}{\usebox{\rstarscaja}}}%
    \fi
    \usebox{\rstarscaja}%
  \else
    \ifdim\wd\rstarscaja>\linewidth
      \resizebox{\linewidth}{!}{\usebox{\rstarscaja}}%
    \else
      \usebox{\rstarscaja}%
    \fi
  \fi}

% (Los colores de los enlaces se fijan en _output.yml: colorlinks, linkcolor...)

% --- Imágenes de capítulo --------------------------------------------
% Las imágenes con pie (las ilustraciones de cada capítulo) no se numeran,
% igual que en la web: solo se muestra el pie, en cursiva y en gris.
\usepackage{caption}
\DeclareCaptionFont{rsgris}{\color{rstarsgris}}
\captionsetup[figure]{labelformat=empty, labelsep=none, font={small,rsgris},
                      skip=6pt}
\setlength{\intextsep}{14pt plus 4pt minus 2pt}

% --- Vídeos -----------------------------------------------------------
% En la web son vídeos de YouTube incrustados; en papel, un recuadro con
% el título, la dirección y un código QR para verlo desde el móvil.
\usepackage{qrcode}
\usepackage[most]{tcolorbox}
\newcommand{\rsvideo}[2]{%
  \par\medskip\noindent
  \begin{tcolorbox}[enhanced, breakable=false, colback=rstarsclaro,
      colframe=rstarsclaro, boxrule=0pt, arc=2mm,
      left=3mm, right=3mm, top=2.5mm, bottom=2.5mm]
    \begin{minipage}[c]{2.3cm}\qrcode[height=2.1cm, nolink]{#2}\end{minipage}%
    \hspace{4mm}%
    \begin{minipage}[c]{\dimexpr\linewidth-2.3cm-4mm\relax}
      {\sffamily\scriptsize\bfseries\color{rstarsgris}\MakeUppercase{Vídeo}}\par
      \smallskip
      {\sffamily\bfseries\color{rstars}#1}\par
      \smallskip
      {\small\href{#2}{\nolinkurl{#2}}}
    \end{minipage}
  \end{tcolorbox}\par\medskip}

% --- Fichas de las compañías (capítulo 6) ------------------------------
% pdf/pdf_hooks.R escribe cada ficha con estos comandos.
\newcommand{\rsbadge}[1]{\fcolorbox{rstarsgris!60}{rstarsclaro}{%
  \rule[-0.25em]{0pt}{1em}\scriptsize\sffamily\bfseries\MakeUppercase{#1}}}
\newtcolorbox{rsficha}{enhanced, breakable=false, colback=white,
  colframe=rstarsgris!45, boxrule=0.6pt, arc=2mm,
  left=3mm, right=3mm, top=2.5mm, bottom=2mm, before skip=8pt, after skip=12pt}
\newcommand{\rsdato}[2]{%
  {\sffamily\tiny\color{rstarsgris}\MakeUppercase{#1}}\par
  {\sffamily\small\bfseries #2}}

% --- Catálogo de planetas (capítulo 4) ---------------------------------
\newcommand{\rsgalaxia}[2]{%
  \par\Needspace{12\baselineskip}\medskip\noindent
  {\sffamily\large\bfseries\color{rstars}#1}%
  \hspace{0.6em}{\sffamily\small\color{rstarsgris}#2}\par
  \nopagebreak\vspace{2pt}{\color{rstarsgris!50}\hrule height 1pt}\par\nopagebreak}
\newcommand{\rsastillero}[1]{\colorbox{rstarsclaro}{%
  \parbox{\dimexpr\linewidth-2\fboxsep\relax}{\raggedright
    \footnotesize\sffamily\bfseries\color{rstars}\strut #1\strut}}}
\usepackage{array}

% --- Contraportada: el libro termina siempre en página par ----------------
% Si la última página con contenido es impar, se añade una en blanco.
\AtEndDocument{%
  \clearpage
  \ifodd\ReadonlyShipoutCounter
    \thispagestyle{empty}\null\clearpage
  \fi}

% --- La portada la pone pdf/portada.tex ---------------------------------
\renewcommand{\maketitle}{}
